\section{The explicit spectral families on wire networks}
\label{sec:explicit-spectrum}
This section computes the wire metrics of the standard explicit spectral points, including the points obtained by grouping parties, and of the more recent laminar extensions.  The distinction between these families and the whole asymptotic spectrum is essential.

Let $V$ be a finite set with at least two elements.  A split of $V$ is an unordered bipartition $S\mid V\setminus S$ into nonempty parts.  Its cut pseudometric is
\[
 \delta_S(u,v)=\begin{cases}1&\text{if exactly one of $u,v$ belongs to $S$,}\\0&\text{otherwise.}\end{cases}
\]
For a loopless multigraph $G$ on $V$, with an integer $q_e\geq1$ assigned to each edge, write
\[
 W_G=\mathop{\bigotimes}_{e=uv\in E(G)}\langle q_e\rangle_{uv},
 \qquad
 r_S(G)=\prod_{\substack{e=uv\in E(G)\\\delta_S(u,v)=1}}q_e.
\]
The tensor products here are taken party by party.  Parallel edges may be combined by multiplying their dimensions. We write $F_\theta=2^{E_\theta}$ for the lower quantum functional and its logarithm; the entropy definition is recalled in the proof below.

\begin{proposition}[Flattenings and entropy on wire products]
\label{prop:wire-evaluations}
Over any field, the matrix flattening of $W_G$ across a split $S\mid V\setminus S$ has rank $r_S(G)$.  Over $\mathbb C$, if $\theta$ is a probability distribution on the splits, the lower quantum functional satisfies
\begin{equation}
 E_\theta(W_G)=\sum_S\theta_S\log_2 r_S(G),
 \qquad
 F_\theta(W_G)=\prod_{e=uv}q_e^{\sum_S\theta_S\delta_S(u,v)}.
 \label{eq:wire-cut-entropy}
\end{equation}
In particular, for singleton weights $\theta_v\geq0$ with $\sum_v\theta_v=1$,
\begin{equation}
 F_\theta(W_G)=\prod_{e=uv}q_e^{\theta_u+\theta_v}.
 \label{eq:wire-star}
\end{equation}
Equation~\eqref{eq:wire-cut-entropy} evaluates a functional.  It does not assert that the functional for an arbitrary split distribution is a spectral point.
\end{proposition}

\begin{proof}
Choose the usual bases on the wire spaces.  A wire crossing the split contributes an identity matrix of rank $q_e$ to the flattening.  A wire wholly on either side contributes a nonzero vector on that side and hence a factor of rank one.  Matrix rank is multiplicative under a Kronecker product, proving the first assertion over every field.

For the entropy assertion, normalize each wire to
$q_e^{-1/2}\sum_{a=1}^{q_e}e_a\otimes e_a$.  The normalized tensor $W_G$ is their tensor product.  Its marginal on $S$ is a tensor product of pure states from internal wires and maximally mixed states from the crossing wires.  Its nonzero spectrum is therefore uniform on $r_S(G)$ entries, with entropy $\log_2 r_S(G)$.

The lower quantum functional is, by Christandl--Vrana--Zuiddam's Definition~3.16 in Section~3.2 of the arXiv version~\cite{cvz-universal}, the supremum of the weighted marginal entropies after invertible maps at the individual parties.  Such maps preserve every flattening rank.  A density matrix of rank $r$ has entropy at most $\log_2 r$: if its positive eigenvalues are $p_1,\ldots,p_r$, this follows from
$\sum_i p_i\log_2(p_i/(1/r))\geq0$.
Consequently every summand in the defining entropy expression is at most $\log_2 r_S(G)$.  The normalized wire tensor simultaneously attains every one of these bounds.  This proves~\eqref{eq:wire-cut-entropy}.  Each edge $uv$ crosses precisely the singleton splits at $u$ and $v$, giving~\eqref{eq:wire-star}.  For two parties the two singleton weights refer to the same split; their sum is one, so the formula still applies.
\end{proof}

The singleton functional is a universal spectral point by Christandl--Vrana--Zuiddam, Theorem~3.30 and Corollary~3.31; its moment-polytope expression is their equation~(25), p.~27 of the arXiv version.  The proposition identifies a simultaneous maximizing spectrum, so no uncomputed moment polytope is needed for these evaluations.  The distinction in its last sentence is substantive: Botero--Christandl--Fraser--Leigh--Nieuwboer prove in Theorem~3.14 of~\cite{bcfln-quantum} that, for at least four parties, every nonsingleton-supported lower quantum functional fails to be a universal spectral point.  Grouping parties before applying a singleton functional is a different operation.

\begin{proposition}[Support functionals on wire products]
\label{prop:wire-support}
For every field and every probability vector $\theta$ on $V$, both Strassen support functionals satisfy
\[
 \zeta^\theta(W_G)=\zeta_\theta(W_G)
     =\prod_{e=uv}q_e^{\theta_u+\theta_v}.
\]
\end{proposition}
\begin{proof}
We use the upper and lower support definitions in Sections~2.1--2.2 of~\cite{cvz-universal}: the logarithmic upper functional minimizes, over bases, the maximum of the weighted marginal Shannon entropies on the support; the lower one maximizes the same entropy quantity on the maximal elements of the support.  We also use their Theorem~2.15, namely $\rho_\theta\leq\rho^\theta$.

Give each edge an orientation and give distinct edges distinct nonnegative integers $j(e)$.  Choose an integer $B$ larger than twice the greatest edge dimension.  At a vertex $v$, for the local tuple $(a_e)_{e\ni v}$ with $0\leq a_e<q_e$, put
\[
 \alpha_v((a_e))=\sum_{e\ni v}\sigma(v,e)B^{j(e)}a_e,
\]
where $\sigma(v,e)$ is $1$ at one endpoint of $e$ and $-1$ at the other.  This is injective: in a nonzero difference of two local tuples, its term of largest exponent has absolute value at least $B^j$, whereas the sum of all smaller terms has absolute value less than $B^j$ by the choice of $B$.
Order the local basis by increasing $\alpha_v$.  On every support tuple of $W_G$, the sum of these labels over the vertices is zero.  Two different support tuples consequently cannot be comparable in the product order.  Thus the support is an antichain and every support element is maximal.

Choose all edge labels independently and uniformly.  At vertex $v$ the marginal is uniform on $\prod_{e\ni v}q_e$ local labels.  The lower functional is therefore at least
$\sum_v\theta_v\sum_{e\ni v}\log_2 q_e$ in logarithmic form.  The upper functional is at most the same number, since no local marginal entropy exceeds the logarithm of the local dimension, and the standard support is one of its permitted choices of basis.  The inequality $\rho_\theta\leq\rho^\theta$ makes both bounds equalities.  Exponentiating proves the claim.
\end{proof}

For complex tensors the upper support functional is now known to equal the singleton quantum functional on all tensors, by Sakabe--Do\u gan--Walter, Theorem~I.1 of~\cite{sdw-support}.  Over arbitrary fields, Alman--Li--Pratt prove universality of the three-party edge support functionals in Theorem~1.1 of~\cite{alp-edge}.  These developments extend the scope of the global spectral points; their values on wire products remain the values in Proposition~\ref{prop:wire-support}.

\subsection{Grouping parties and the closed family of clustered stars}
The next result includes all singleton quantum points obtained after grouping several of the original parties.

Let $\pi$ be a partition of $V$ into at least two nonempty blocks.  Regard a $V$-tensor as a tensor whose parties are the blocks, and apply a singleton quantum functional with probability weights $\theta_B$ on those blocks.  Grouping respects products, direct sums up to zero ambient coordinates, and restrictions.  This composition is therefore a universal spectral point.  Denote the collection of these points, for all $\pi$, by $\mathcal Q_V$.

\begin{theorem}[The scaled metrics of grouped singleton functionals]
\label{thm:grouped-stars}
The scaled wire metrics of $\mathcal Q_V$ form exactly the following set:
choose a partition $\pi$ of $V$ into at least two blocks and nonnegative numbers $a_B$, and set
\begin{equation}
 d(u,v)=\begin{cases}
 0,&B(u)=B(v),\\
 a_{B(u)}+a_{B(v)},&B(u)\ne B(v).
 \end{cases}
 \label{eq:clustered-star}
\end{equation}
This is a closed finite union of polyhedral cones.  For three labels it is the whole pseudometric cone.
\end{theorem}
\begin{proof}
A wire whose endpoints are in the same block becomes a nonzero simple tensor in the grouped system and has value one.  A wire joining different blocks becomes a bipartite wire between those blocks.  Proposition~\ref{prop:wire-evaluations} gives exponent $\theta_{B(u)}+\theta_{B(v)}$.  Scaling by $c\geq0$ gives~\eqref{eq:clustered-star} with $a_B=c\theta_B$.
Conversely, when $c=\sum_Ba_B>0$, choose $\theta_B=a_B/c$.  When all $a_B$ vanish the zero scale realizes the zero metric.

For a fixed partition the cone is the nonnegative span of its finitely many block-cut metrics $\delta_B$ and is therefore polyhedral and closed.  There are only finitely many partitions of $V$.  Finally, for three labels $1,2,3$, every pseudometric is a star with
\[
 a_1=\frac{d_{12}+d_{13}-d_{23}}2,\quad
 a_2=\frac{d_{12}+d_{23}-d_{13}}2,\quad
 a_3=\frac{d_{13}+d_{23}-d_{12}}2.
\]
The triangle inequalities make these numbers nonnegative, and substitution gives the prescribed distances.
\end{proof}

\subsection{Laminar extensions and tree metrics}
The 2026 laminar extension theorem supplies more global spectral points, although it gives their extension by existence rather than by an explicit formula on every tensor.  Its wire values admit a complete geometric description.

Two splits are \emph{compatible} if at least one of their four pairwise intersections is empty.  A split distribution is \emph{laminar} if all splits with positive weight are pairwise compatible.  Botero--Christandl--Fraser--Leigh--Nieuwboer's Proposition~4.3~\cite{bcfln-quantum}, using Strassen's extension theorem, extends the common upper and lower laminar quantum functional from their semiring of laminar-quantum-robust tensors to a universal spectral point.  Wire products belong to this semiring, by their Proposition~4.12 and closure under products.  Let $\mathcal L_V$ be the collection of all these extensions.

\begin{theorem}[The wire metrics of laminar extensions]
\label{thm:laminar-tree}
The scaled wire metrics of $\mathcal L_V$ are exactly the pseudometrics obtained by mapping $V$ into a finite tree with nonnegative edge lengths.  Their set is closed.  For three labels it is one cone with three pendant-edge generators; for four labels it is the union of the three binary-tree cones, each with five edge generators; for five labels it is the union of the fifteen binary-tree cones, each with seven edge generators.  Zero edge lengths include all degeneracies.
\end{theorem}
\begin{proof}
By Proposition~\ref{prop:wire-evaluations}, an extension associated with $\theta$ has wire metric $\sum_S\theta_S\delta_S$.  Scaling removes the normalization $\sum_S\theta_S=1$.  We show directly that the nonnegative sums of compatible cut metrics are precisely the tree pseudometrics.

Deleting any edge of a tree separates its labels into two parts, and the distance between two labels is the sum of the lengths of the edges separating them.  The resulting edge splits are pairwise compatible.  Edges with no label on one side may be removed, and labels at internal vertices may be given pendant edges of length zero.

Conversely, choose a distinguished label $r\in V$.  Represent each compatible split by the side not containing $r$.  These subsets form a laminar collection: any two are disjoint, or one contains the other.  Add $V$ and all singletons, assigning zero length to subsets not already present.  Make a vertex for each of these subsets.  Join each proper subset to the smallest strictly larger subset in the collection that contains it.  The parent is unique, since two candidate parents intersect and are therefore nested.  The resulting connected graph is a tree: every vertex except $V$ has one parent, and inclusion strictly increases along every parent edge.  Give the edge from $S$ to its parent the coefficient of $\delta_S$, or zero when $S$ was added.  Put label $v$ at vertex $\{v\}$.  The edge from $S$ to its parent separates precisely the labels in $S$ from the others, so its distance function is the required cut sum.

There are finitely many possible compatible split families on a fixed finite label set, and each determines a finitely generated closed cone.  Their union is closed.  Any nonbinary tree may be resolved into a binary tree by adding edges of length zero.  There is one binary unrooted shape with three labelled leaves, three with four, and fifteen with five: a four-leaf binary tree is determined by its $2\mid2$ split; a five-leaf binary tree is determined by choosing the middle leaf in five ways and partitioning the other four labels into its two unordered cherries in three ways.  A binary tree with $n$ leaves has $2n-3$ edges, by the degree sum and Euler's formula.  This gives the stated counts of generators.
\end{proof}

\begin{theorem}[Exact distance from the square]
\label{thm:square-tree-gap}
Let $e$ be the metric of a unit square, labelled cyclically $1,2,3,4$, and let $e'$ include its centre $0$.  In the supremum norm on unordered pairs,
\[
 \operatorname{dist}(e,\{c\,d_F:F\in\mathcal Q_V,\ c\geq0\})
 =\operatorname{dist}(e,\{c\,d_F:F\in\mathcal L_V,\ c\geq0\})
 =\frac{\sqrt2-1}{2}.
\]
The same two equalities hold for $e'$ on its five labels.
\end{theorem}
\begin{proof}
For any four points in a tree, the two largest of
\[
 d_{12}+d_{34},\qquad d_{13}+d_{24},\qquad d_{14}+d_{23}
\]
are equal.  Indeed, after pruning and suppressing unlabelled degree-two vertices, the minimal subtree connecting the points has either one central vertex, in which case all three sums agree, or a central path separating two pairs.  Each sum contains every pendant segment once.  The pairing within the two sides omits the central path, and each of the other pairings contains it twice.  Zero lengths and coincident labels follow by the same formula.

For the square the three sums are $2,2\sqrt2,2$.  If a tree pseudometric $d$ is within $\varepsilon$ of $e$, its middle-labelled sum is at least $2\sqrt2-2\varepsilon$, while the other two sums are at most $2+2\varepsilon$.  If $\varepsilon<(\sqrt2-1)/2$, the middle-labelled sum would be strictly larger than both others, contradicting the tree property.  Restricting a tree metric on the square and centre to the four corners gives the same lower bound for $e'$.

For the matching upper bound, put four labelled leaves on a star, each with pendant length
\[
 a=\frac{1+\sqrt2}{4}.
\]
Every distance between corners is $2a=(1+\sqrt2)/2$, whose deviations from $1$ and $\sqrt2$ are both $(\sqrt2-1)/2$.  For $e'$, add a fifth leaf labelled $0$, of pendant length
\[
 a_0=\frac{\sqrt2-1}{4}.
\]
Then its distance to each corner is $a_0+a=\sqrt2/2$, exactly the required radial distance.  These stars belong to $\mathcal Q_V$ after scale by Theorem~\ref{thm:grouped-stars}; they also belong to $\mathcal L_V$.  Finally every clustered star is itself a tree pseudometric, so the lower bound applies to both families.
\end{proof}

Theorem~\ref{thm:square-tree-gap} excludes these spectral constructions, including their pointwise limits.  It does not exclude the square from the price cone.  In particular it is not a classification of all spectral points, and the existence of other higher-party spectral points is already proved in~\cite{alp-edge,bcfln-quantum}.

\begin{corollary}[Cross-polytope configurations]
\label{cor:cross-polytope-tree-gap}
For $D\geq2$, give $X_D=\{\pm e_1,\ldots,\pm e_D\}$ its Euclidean metric, with or without the centre $0$.  Its exact distance in the sup norm to both the scaled grouped-singleton family and the scaled laminar-tree family is
\[
 1-\frac1{\sqrt2}.
\]
For $D=1$ the distance is zero.
\end{corollary}
\begin{proof}
Two antipodal pairs form a square with side $\sqrt2$ and diagonals two.  The four-point argument in Theorem~\ref{thm:square-tree-gap} therefore gives the lower bound $\sqrt2(\sqrt2-1)/2=1-1/\sqrt2$.  For the upper bound take a star whose terminal pendant edges all have length $(2+\sqrt2)/4$.  Every terminal distance is $(2+\sqrt2)/2$, at exactly the asserted error from either possible Euclidean distance, $\sqrt2$ or two.  If the centre is included, give it a pendant edge of length $(2-\sqrt2)/4$; all its terminal distances are then exactly one.  These are singleton star metrics.  The configurations for $D=1$ lie on a line and are already tree metrics.
\end{proof}

\begin{corollary}[Extra parties do not remove this gap]
\label{cor:tree-pullbacks}
Starting with any of the tree metrics of Theorem~\ref{thm:laminar-tree}, adding finitely many parties, pulling back along maps of label sets, and taking pointwise limits cannot produce the square or the square with its centre.
\end{corollary}
\begin{proof}
A pullback is obtained by composing the map into the same weighted tree with a map of label sets.  It is again a tree pseudometric.  The tree-pseudometric cone on the final finite label set is closed, and Theorem~\ref{thm:square-tree-gap} gives a positive distance from it.
\end{proof}

\subsection{An exact formal candidate and its missing extension}
The following calculation records why excluding trees alone leaves the main problem open.
Put $a=1-\sqrt2/2$ and $b=\sqrt2-1$.  On the four corners,
\begin{equation}
 e=a\sum_{i=1}^{4}\delta_{\{i\}}
       +b\bigl(\delta_{\{1,2\}}+\delta_{\{1,4\}}\bigr).
 \label{eq:square-cut-decomposition}
\end{equation}
A side crosses one of the two displayed $2\mid2$ cuts and two singleton cuts, so its value is $2a+b=1$.  A diagonal crosses both of those cuts and two singleton cuts, so its value is $2a+2b=\sqrt2$.  The total cut weight is $4a+2b=2$.  The two nonsingleton splits are incompatible.
For the square and centre, set $t=(\sqrt2-1)/2$.  Then
\begin{equation}
 e'=a\sum_{i=1}^{4}\delta_{\{i\}}
       +t\bigl(\delta_{\{1,2\}}+\delta_{\{2,3\}}
                 +\delta_{\{3,4\}}+\delta_{\{4,1\}}\bigr).
 \label{eq:centred-square-cut-decomposition}
\end{equation}
Here a side has value $2a+2t=1$, a diagonal has value $2a+4t=\sqrt2$, and a radial pair has value $a+2t=\sqrt2/2$.  Again the weights sum to two.

Thus the lower quantum functionals for half of the weights in~\eqref{eq:square-cut-decomposition} and~\eqref{eq:centred-square-cut-decomposition} have exactly the wire metrics $e/2$ and $e'/2$.  Theorem~3.14 of~\cite{bcfln-quantum} rules out using these lower functionals themselves as universal spectral points.  A universal spectral point agreeing with them on wire products, if one exists, must be obtained by a further extension argument.  The laminar extension theorem does not supply that argument, because the required cuts cross.  This is a precise obstruction to the proposed entropy construction, not a proof of nonexistence of the desired spectral point.

\subsection{Why singleton-anchor interpolation cannot create a flat square}
The obstruction also applies to interpolation from an arbitrary price, not merely from the families classified above.  This uses the quadrilateral inequality for prices proved through the CAT$(0)$ embedding theorem.

\begin{proposition}[The square has no positive pendant summand]
\label{prop:square-no-pendant}
Assume that every price $d$ satisfies
\begin{equation}
 d_{13}^2+d_{24}^2
 \leq d_{12}^2+d_{23}^2+d_{34}^2+d_{41}^2.
 \label{eq:price-quadrilateral}
\end{equation}
If $a_1,\ldots,a_4\geq0$ and the off-diagonal function
$d_{ij}=e_{ij}-a_i-a_j$ is a price, then every $a_i$ is zero.  For the square with centre, the same conclusion holds for $a_0,\ldots,a_4\geq0$ and $d_{ij}=e'_{ij}-a_i-a_j$.
\end{proposition}
\begin{proof}
Put $t=1-\sqrt2/2>0$ and $s=a_1+a_2+a_3+a_4$.  The triangle inequality for the triple consisting of $i$ and its two cyclic neighbours gives
\[
 0\leq d_{i-1,i}+d_{i,i+1}-d_{i-1,i+1}
       =2-\sqrt2-2a_i.
\]
Thus $0\leq a_i\leq t$ for every corner.  Let $D(d)$ be the left side of~\eqref{eq:price-quadrilateral} minus its right side.  Direct expansion gives
\[
 D(d)=4ts-\sum_i a_i^2
       +2(a_1a_3+a_2a_4)
       -2(a_1a_2+a_2a_3+a_3a_4+a_4a_1).
\]
Since $a_i^2\leq ta_i$ and $2a_ia_j\leq t(a_i+a_j)$, the two negative sums have total absolute value at most $ts+2ts=3ts$.  The opposite-pair contribution is nonnegative.  Hence $D(d)\geq ts$.  The assumed quadrilateral inequality forces $s=0$ and therefore every corner coefficient vanishes.

For five labels, restriction to the corners gives the same conclusion.  In fact, regardless of the corner coefficients, the triangle inequality along either diameter through the centre gives
\[
 0\leq d_{10}+d_{03}-d_{13}=-2a_0,
\]
so $a_0=0$ as well.
\end{proof}

Wigderson--Zuiddam's graph-spectrum interpolation theorem is log-star-convexity toward a single-party flattening; see Section~12.7, Theorem~12.23 in the author version of~\cite{wz-spectra}.  In wire exponents this is interpolation $d\mapsto(1-\lambda)d+\lambda\delta_{\{v\}}$.  Proposition~\ref{prop:square-no-pendant} shows that no finite succession of such interpolations with a positive final pendant contribution can give either square configuration, even if the starting price is arbitrary.  Local changes of tensor bases give no additional freedom: restriction monotonicity in both directions makes every spectral point invariant under invertible maps at the parties.  The symmetric-space and moment-polytope descriptions may still lead to other constructions, but neither coordinate changes nor these known anchor interpolations supply one.

\subsection{Every crossing four-party upper quantum expression fails}
The formal cut decomposition does not become a spectral point by replacing the lower quantum functional with the ordered upper quantum functional.  In fact every four-party weighting with crossing split support fails already on the rank-two unit tensor.  This strengthens the crossing-projector example in Appendix~D of Botero--Christandl--Fraser--Leigh--Nieuwboer~\cite{bcfln-quantum} by allowing arbitrary singleton weights and all three nonsingleton splits.

For clarity, the ordered upper expression uses a different definition.
For a partition $\lambda\vdash n$, let $H(\lambda/n)$ be the base-two
Shannon entropy of its normalized row lengths, and let
$P^{(b)}_\lambda$ be the Schur--Weyl isotypic projector for the grouped
cut $b$, on the globally symmetric subspace of the $n$ tensor copies.
For a nonzero tensor $T$ and an ordering $b_1,\ldots,b_s$ of the
positive-weight cuts, the definition in
\cite[Definition 3.3]{cvz-universal} is
\[
 E^\theta(T)=\sup\left\{
   \sum_{j=1}^s\theta(b_j)H(\lambda_j/n):
   n\geq1,\ \lambda_j\vdash n,\quad
   P^{(b_s)}_{\lambda_s}\cdots P^{(b_1)}_{\lambda_1}
   T^{\otimes n}\ne0\right\}.
\]
We put $F^\theta(T)=2^{E^\theta(T)}$ and $F^\theta(0)=0$.
The order is part of the expression when the cut projectors do not
commute. Thus any one nonzero ordered projector tuple proves the
corresponding entropy lower bound.

\begin{proposition}[Failure of crossing upper quantum functionals]
\label{prop:square-upper-fails}
Let $\theta$ be a probability distribution on the bipartitions of four parties, with at least two positive weights on the three $2\mid2$ splits.  List those positive-weight splits in the order in which their projectors are applied, as $b_1,\ldots,b_r$, where $r\in\{2,3\}$.  For the ordered upper quantum functional and the unnormalised unit tensor $U_2=\sum_{i=0}^{1}e_i^{\otimes4}$,
\begin{equation}
 E^\theta(U_2)\geq 1+\sum_{j=2}^r\theta(b_j)>1.
 \label{eq:upper-square-failure}
\end{equation}
Thus $F^\theta$ fails even the cover inequality for relaxed spectral points.  For the square weights in~\eqref{eq:square-cut-decomposition}, divided by two, either order gives
\[
 E^\theta(U_2)\geq\frac{1+\sqrt2}{2}>1.
\]
\end{proposition}
\begin{proof}
We use four tensor copies.  Assign Young shape $(2,2)$ to every singleton split and to $b_1$, and shape $(1,1,1,1)$ to every subsequent $2\mid2$ split.  We prove that the corresponding ordered projector product on $U_2^{\otimes4}$ is nonzero.  The normalized shapes have entropies $1$ and $2$, respectively; nonvanishing therefore gives precisely~\eqref{eq:upper-square-failure}.

Singleton projectors commute with all cut projectors.  At one party the weight-two subspace of $(\mathbb C^2)^{\otimes4}$ has orthonormal basis $e_X$, indexed by the two-element subsets of $\{1,2,3,4\}$.  Choose
\[
 X_1=\{1,2\},\qquad X_2=\{1,3\},\qquad X_3=\{1,4\},
\]
one representative from each complementary pair, and put
\[
 u_r=\frac{e_{X_r}+e_{X_r^c}}{\sqrt2},\qquad
 q_r=u_r-\frac{u_1+u_2+u_3}{3}.
\]
The span $E$ of the $q_r$ is the two-dimensional irreducible $S_4$-module $[2,2]$.  Indeed the complementary-pair sums form the permutation representation on three objects, whose nonconstant subspace is $[2,2]$.  All other Hamming-weight sectors have only $[4]$ and $[3,1]$, so the singleton $[2,2]$ projectors kill them.  Since $P_{22}e_{X_r}=P_{22}e_{X_r^c}=q_r/\sqrt2$, the product of all four singleton projectors sends $U_2^{\otimes4}$ to
\[
 T_0=\frac12\sum_{r=1}^3q_r^{\otimes4}.
\]
The vectors $q_r$ have squared norm $2/3$ and pairwise inner product $-1/3$.  In an orthonormal basis $f_1,f_2$ of $E$ they can be taken as
\[
 q_1=(\sqrt{2/3},0),\quad
 q_2=(-1/\sqrt6,1/\sqrt2),\quad
 q_3=(-1/\sqrt6,-1/\sqrt2).
\]
A direct expansion gives
\begin{equation}
 T_0=\frac1{12}(I_1+I_2+I_3),
 \label{eq:fourth-moment-pairings}
\end{equation}
where the three pairing tensors have coordinates
\[
 (I_1)_{abcd}=\delta_{ab}\delta_{cd},\qquad
 (I_2)_{abcd}=\delta_{ac}\delta_{bd},\qquad
 (I_3)_{abcd}=\delta_{ad}\delta_{bc}.
\]
Here each of $a,b,c,d$ ranges over $\{1,2\}$.

On $E\otimes E$, the decomposition is
\[
 [2,2]\otimes[2,2]=[4]\oplus[2,2]\oplus[1,1,1,1].
\]
The three summands are respectively the scalar symmetric tensors, the traceless symmetric tensors, and the antisymmetric tensors.  This follows either from the $S_4$ characters or from the standard two-dimensional representation of $S_3$ through which $[2,2]$ factors.  Consequently the $22$ projector takes the symmetric part and subtracts its trace, while the $1111$ projector takes the antisymmetric part.

For two crossing cuts we may relabel the parties so that $b_1=AD\mid BC$ and $b_2=AB\mid CD$.  Represent the first projector on $BC$ and the second on $AB$.  Applying the traceless-symmetric $BC$ projector to~\eqref{eq:fourth-moment-pairings} gives
\[
 T_1=\frac1{12}(I_1+I_2-I_3).
\]
Antisymmetrizing its $AB$ factors then gives
\[
 T_2=\frac1{12}(I_2-I_3)\ne0.
\]
For three cuts, every ordering can similarly be relabeled to $BC\mid AD$, $AC\mid BD$, $AB\mid CD$.  After the same first step, antisymmetrizing $AC$ gives $(I_1-I_3)/12$, and antisymmetrizing $AB$ gives
\[
 T_3=\frac1{24}(I_2-I_3)\ne0.
\]
Party relabelings induce every permutation of the three $2\mid2$ splits, covering all orders.

These calculations also cover zero singleton weights.  If a singleton cut has zero weight, its projector is absent from the definition.  A nonzero result after applying the additional commuting singleton projectors implies nonvanishing before those additional projections.  All intermediate tensors are invariant under simultaneous permutations of the four copies; hence the total symmetrizers in Definition~3.3 of~\cite{cvz-universal} make no further change.

Each simple nonzero tensor has upper functional one, since its surviving Young shapes have only one row.  But $U_2$ is the sum of two such tensors, whereas~\eqref{eq:upper-square-failure} gives $F^\theta(U_2)>2$.  This violates the cover inequality.  The square has both nonsingleton weights $(\sqrt2-1)/2$, giving the displayed special case.
\end{proof}

For additional exact checks, $\|I_j\|^2=4$ and $\langle I_i,I_j\rangle=2$ for $i\ne j$.  Thus $\|T_2\|^2=1/36$ and $\|T_3\|^2=1/144$.  At the basis coordinate $(X_1,X_3,X_3,X_1)$ their coefficients are $-1/144$ and $-1/288$.  The exact character-projector program \emph{Quantum crossing projectors} in the \href{\workbenchurl}{workbench} reproduces these values, checks every cut ordering, and includes a second contraction using the three pairing tensors.  This excludes the indicated upper expressions on all tensors; a different universal point could still agree with the formal cut expression on wire products.

\begin{corollary}[The centred-square upper expression also fails]
\label{cor:centred-square-upper-fails}
Let $\theta$ be half of the cut weights in~\eqref{eq:centred-square-cut-decomposition}.  For every ordering of its four adjacent-pair cut projectors, the upper quantum functional on five parties satisfies
\[
 E^\theta(U_2(\{0,1,2,3,4\}))\geq\frac{1+\sqrt2}{2}>1.
\]
In particular this upper expression is not a relaxed spectral point.
\end{corollary}
\begin{proof}
Use four copies, shape $22$ at the four corners, and an additional shape $4$ projector at the centre, whose weight is zero.  On the weight-two subspace the centre projector takes every basis word to
\[
 s=\frac16\sum_{|X|=2}e_X\ne0.
\]
The same calculation as in the preceding proof therefore sends the five-party unit tensor's fourth power to $T_0\otimes s$, with $T_0$ given by~\eqref{eq:fourth-moment-pairings}.  The extra centre projector commutes with every listed projector, since the adjacent-pair representatives contain only corners.  It is enough to prove nonvanishing after this additional projection.

On the globally invariant four-corner subspace, complementary cuts have equal isotypic projectors.  Thus the four adjacent-pair cuts act as two families: $12$ and $34$ give one family, and $23$ and $41$ the other.  Assign shape $22$ to the initial consecutive block in the first family encountered, and shape $1111$ to every later cut.  The initial block has length one or two and its repeated projector is idempotent.  At the first change of family the preceding proof gives a nonzero antisymmetric pairing tensor.  Subsequent $1111$ projectors cannot kill it.  Indeed the antisymmetric lines for the two families are spanned, up to sign and permutation, by $I_2-I_3$ and $I_1-I_2$; their squared norms are four and their mutual inner product is $-2$.  Projection from either line to the other multiplies by a nonzero scalar, while projection onto the same line leaves it unchanged.  This proves nonvanishing for every cut order.

Write $w=(\sqrt2-1)/4$ for each of the four adjacent-pair weights.  Every positive-weight cut has entropy at least one in this witness, and at least two of the adjacent-pair cuts have entropy two.  Since the weights sum to one, the resulting entropy is at least $1+2w=(1+\sqrt2)/2$.  The rank-two unit tensor is a sum of two simple tensors, so its upper functional again violates the cover inequality.
\end{proof}

\subsection{Tensor-power regularization does not repair the crossing expression}

A positive power can reduce an excessive value on a unit tensor.
The preceding cover obstruction alone therefore does not exclude
all powers of the ordered upper expression. Its failure of
multiplicativity is stronger, and persists at every wire dimension.
Here $F^\theta$ always denotes the fixed-order expression from
\cite[Definition~3.3]{cvz-universal}; no submultiplicativity property
is assumed outside that source's noncrossing hypotheses.

\begin{proposition}[Exact upper value on two matching wires]
\label{prop:upper-matching-wires}
Let $\theta$ be a probability distribution on the bipartitions of
four parties. Write $s$ for its total singleton weight and $c$
for its total $2\mid2$ weight. Suppose at least two $2\mid2$
cuts have positive weight, and list them as $b_1,\ldots,b_r$
in the order in which their projectors act, where $r\in\{2,3\}$.
Fix $j>1$. For $q\ge2$, let $W_q=X_qY_q$ be the product of
one dimension-$q$ wire on each pair of the matching $b_j$.
Then
\begin{align}
 F^\theta(W_q)&=q^{s+2c},\label{eq:upper-matching-value}\\
 F^\theta(X_q)F^\theta(Y_q)
   &=q^{s+2(c-\theta(b_j))}<F^\theta(W_q).
 \label{eq:upper-matching-gap}
\end{align}
Arbitrary singleton weights and every order of their commuting
projectors are allowed.
\end{proposition}
\begin{proof}
Each local space of $W_q$ has dimension $q$. A singleton Young
shape therefore has at most $q$ rows, and a two-party shape
has at most $q^2$ rows. The ambient entropy bound gives
$E^\theta(W_q)\le(s+2c)\log_2q$. This uses dimensions, not
flattening ranks after noncommuting projections.

For equality take $n=q^2$ copies. Assign the square partition
$\lambda=(q^q)$ to every singleton and the sign partition
$(1^n)$ to every $2\mid2$ cut. Since $b_1\ne b_j$, the tensor
$W_q$ is maximally entangled of dimension $q^2$ across $b_1$.
The first sign projector is consequently nonzero on $W_q^{\otimes n}$
and leaves a product of the two top exterior vectors across that cut.
The representation
\[
 \bigwedge^{q^2}(\mathbb C^q\otimes\mathbb C^q)
\]
transforms as $\det^q$ on each factor. Thus the singleton Schur
shape is $\lambda$ at every party, its general-linear-group
factor is one-dimensional, and the remaining symmetric-group
factor is $E=[\lambda]$.

The square partition is self-conjugate. Hence the sign
representation occurs with multiplicity one in $E\otimes E$.
Choose a real orthogonal realization of $E$ and let $J$ represent
this sign pairing. Schur's lemma permits the normalization
\[
 J^{\mathsf T}J=I,\qquad J^{\mathsf T}=\varepsilon J,
 \qquad J^2=\varepsilon I,\qquad \varepsilon\in\{1,-1\}.
\]
For completeness, a nonzero real sign intertwiner exists because
$[\lambda]\otimes\operatorname{sgn}\cong[\lambda]$.
Its positive operator $J^{\mathsf T}J$ is scalar; after normalizing
it, $J^2$ is scalar and orthogonal, giving the displayed signs.
Put $d=\dim E$.

On the globally permutation-invariant part of $E^{\otimes4}$,
the sign range for a pair cut is the line of its two $J$-pairings.
The three lines have representatives
\[
 J_{ab}J_{cd},\qquad J_{ac}J_{bd},\qquad J_{ad}J_{bc}.
\]
Each has squared norm $d^2$, and every two distinct representatives
have inner product $d$ or $\varepsilon d$, up to orientation signs.
For example,
\[
 \sum_{a,b,c,d}J_{ab}J_{cd}J_{ad}J_{bc}
   =\sum_{b,c}J_{cb}J_{bc}=\operatorname{tr}(J^2)
   =\varepsilon d\ne0.
\]
The other overlaps follow by exchanging indices. Therefore every
later sign projector retains a nonzero vector. The singleton
projectors commute with all cut projectors and fix these vectors,
so they can be inserted in their prescribed positions. The
surviving shapes have entropies $\log_2q$ and $2\log_2q$,
which attain the ambient bound and prove
\eqref{eq:upper-matching-value}.

For a single wire, cuts separating its endpoints carry the same
nontrivial Young shape; every other cut has a one-row shape.
Its upper value is consequently $q$ raised to the total weight
of separating cuts. The four singleton weights contribute $s$
to the product of the two wire values. The internal matching cut
contributes zero, and each other $2\mid2$ cut separates both
wires and contributes twice its weight. This proves
\eqref{eq:upper-matching-gap}.
\end{proof}

\begin{corollary}[Positive powers and tensor-power limits still fail]
\label{cor:upper-power-regularization-fails}
For the crossing distribution in
Proposition~\ref{prop:upper-matching-wires}, define
\[
 \overline F(T)=\limsup_{m\to\infty}
                  F^\theta(T^{\otimes m})^{1/m},\qquad
 \underline F(T)=\liminf_{m\to\infty}
                  F^\theta(T^{\otimes m})^{1/m}.
\]
Both are finite, and neither is multiplicative. Every positive
power of $F^\theta$, $\overline F$, or $\underline F$ fails
multiplicativity as well. In particular, if the ordinary
tensor-power limit exists on all tensors, it is not a relaxed
spectral point after any positive power.
\end{corollary}
\begin{proof}
The ambient entropy bound scales linearly with $m$, proving
finiteness. Grouped tensor powers satisfy
$W_q^{\otimes m}\cong W_{q^m}$, with the same identities for
each separate wire. Equations~\eqref{eq:upper-matching-value}
and~\eqref{eq:upper-matching-gap} therefore hold at every power.
On these three tensors, the limsup and liminf equal the original
values exactly. Their strict multiplicativity gap survives every
positive power.
\end{proof}

\begin{corollary}[The same obstruction for the centred square]
\label{cor:centred-upper-power-fails}
Use the centred-square weights of
Corollary~\ref{cor:centred-square-upper-fails}, and give the centre
party a one-dimensional factor. For every order of the four
adjacent-pair cuts, choose $W_q=X_qY_q$ from the matching family
which is not applied first. Then
\[
 F^\theta(W_q)=q^{\sqrt2},\qquad
 F^\theta(X_q)F^\theta(Y_q)=q.
\]
Thus all the tensor-power and positive-power conclusions of
Corollary~\ref{cor:upper-power-regularization-fails} hold for this
five-party expression.
\end{corollary}
\begin{proof}
On globally permutation-invariant corner tensors, the cuts $12$
and $34$ give the same projector, as do $23$ and $41$.
The centre factor changes neither equality. Choose the matching
from the family other than the first one. The proof of
Proposition~\ref{prop:upper-matching-wires} applies with sign
shapes at every pair cut. Repetitions of a family are idempotent;
switching families has nonzero pairing overlap. The singleton
weights sum to $s=2-\sqrt2$ and the pair weights to $c=\sqrt2-1$,
so the full upper exponent is $s+2c=\sqrt2$.
The two internal-family weights sum to $(\sqrt2-1)/2$,
leaving exponent one in the product of the individual-wire
values. The same identities hold after replacing $q$ by $q^m$.
\end{proof}

These results concern the ordered projector expression and the
specified tensor-power regularizations. A different spectral
extension or envelope might change its values on products of
wires while retaining individual-wire exponents. Such a
construction is not excluded here. In particular, the abstract
regularization theorems for upper functionals require genuine
subadditivity and submultiplicativity as hypotheses
\cite[arXiv version~2, Definitions~3.1 and~3.13]{BugarVrana2025}; the crossing
expression above already fails the latter.

\subsection{Collinear triples force edge support functionals}
The new edge-support classification gives a different necessary condition on a hypothetical square price.  It determines the spectral values of every tensor on either diameter and its midpoint, not merely the three wire values.

\begin{proposition}[Spectral rigidity on a collinear triple]
\label{prop:collinear-edge-point}
Let $\phi$ be a relaxed spectral point over any field on a finite party set $V$.  Suppose distinct parties $x,z,y$ satisfy
\[
 d_\phi(x,y)=d_\phi(x,z)+d_\phi(z,y)=g>0.
\]
For every tensor $T$ supported on $\{x,z,y\}$,
\begin{equation}
 \phi(T)=\bigl(\zeta^{(\rho,1-\rho,0)}(T)\bigr)^g,
 \qquad \rho=\frac{d_\phi(x,z)}{g},
 \label{eq:collinear-edge-value}
\end{equation}
where the three slots on the right are ordered $(x,y,z)$ and $\zeta$ is the edge support functional of Alman--Li--Pratt.
\end{proposition}
\begin{proof}
Put $P=\{x,z,y\}$ and let $g_P$ be defined by $\phi(U_q(P))=q^{g_P}$.  Applying the linear functional that takes every basis vector to one at $z$ restricts $U_q(P)$ to $\langle q\rangle_{xy}$, so $g_P\geq d_\phi(x,y)$.  Conversely $U_q(P)$ is a restriction of
$\langle q\rangle_{xz}\langle q\rangle_{zy}$: at $z$, send $e_i\otimes e_j$ to $e_i$ when $i=j$ and to zero otherwise.  Hence
\[
 d_\phi(x,y)\leq g_P\leq d_\phi(x,z)+d_\phi(z,y).
\]
The collinearity assumption forces $g_P=g$.

Restriction of $\phi$ to tensors supported on $P$ is a relaxed spectral point.  By Theorem~\ref{thm:normalization}, its $1/g$ power is an ordinary normalized spectral point $\eta$ on three parties.  It satisfies
$\eta(\langle q\rangle_{xy})=q$ for every $q\geq1$.
Corollary~8.10 of~\cite[page 53]{alp-edge} therefore makes $\eta$ the edge support functional $\zeta^{(\rho,1-\rho,0)}$ for some $\rho\in[0,1]$, in the party order $(x,y,z)$.  Proposition~\ref{prop:wire-support} identifies its $xz$ wire exponent as $\rho$, giving the value in~\eqref{eq:collinear-edge-value}.
\end{proof}

\begin{corollary}[Two prescribed three-party restrictions]
\label{cor:square-edge-restrictions}
If the square with its centre is a price over $\mathbb C$, there are an ordinary normalized five-party spectral point $\psi$ and $c>0$ with $e'=c\,d_\psi$.  Write $g=\sqrt2/c$.  On each of the triples $\{1,0,3\}$ and $\{2,0,4\}$ its restriction is
\[
 T\longmapsto\bigl(\zeta^{(1/2,1/2,0)}(T)\bigr)^g,
\]
with the centre in the zero-weight slot.  In addition, its four adjacent-corner wire exponents are all $g/\sqrt2$, its two diagonal exponents are $g$, and its four radial exponents are $g/2$.
\end{corollary}
\begin{proof}
Corollary~\ref{cor:price-attained} supplies $c,\psi$.  Each diameter has its centre as an exact midpoint.  Apply Proposition~\ref{prop:collinear-edge-point} to the two diameter triples; the remaining exponents are the prescribed distances divided by $c$.
\end{proof}

The two three-party restrictions alone are compatible.  For example the singleton quantum point with centre weight zero and each corner weight $1/4$ has both of these restrictions with $g=1/2$, but its adjacent and diagonal exponents are equal.  Thus the extension problem in Corollary~\ref{cor:square-edge-restrictions} includes the four additional adjacent-wire values.  Proving or excluding a common extension with adjacent-to-diagonal ratio $1/\sqrt2$ remains a concrete form of the square question.  The corollary establishes necessary spectral values on two full three-party tensor semirings; it does not assert that the prescriptions cannot coexist.

\begin{proposition}[A saturated wire forces a line metric]
\label{prop:saturated-wire-line}
Assume the price-space theorem, namely that every finite price is induced by a map into a CAT$(0)$ space.  Let $\psi$ be an ordinary normalized spectral point on $V$.  If $d_\psi(x,y)=1$ for some pair $x,y$, then
\begin{equation}
 d_\psi(z,w)=\bigl|d_\psi(x,z)-d_\psi(x,w)\bigr|
 \quad\text{for all }z,w\in V.
 \label{eq:saturated-line}
\end{equation}
In particular its wire pseudometric is induced by a map into the unit interval.
\end{proposition}
\begin{proof}
Fix $z\notin\{x,y\}$ and let $g_P$ be the unit-tensor exponent on $P=\{x,y,z\}$.  The restriction of $U_q(P)$ to the $xy$ wire gives $g_P\geq1$.  Contracting the factors outside $P$ restricts $U_q(V)$ to $U_q(P)$, and normalization of $\psi$ gives $g_P\leq1$.  Thus $g_P=1$, and Theorem~\ref{thm:normalization} makes the restriction of $\psi$ to $P$ an ordinary normalized three-party point.  Alman--Li--Pratt's Corollary~8.10 identifies it as $\zeta^{(\rho,1-\rho,0)}$, in the order $(x,y,z)$.  Evaluation on its other two wires yields
\[
 d_\psi(x,z)+d_\psi(z,y)=1=d_\psi(x,y).
\]
In a CAT$(0)$ realization, this triangle equality puts the image of $z$ on the unique geodesic segment from the image of $x$ to that of $y$.  Indeed the concatenation of the two geodesic segments through $z$ has length equal to the distance between the endpoints, hence is a geodesic; geodesic uniqueness identifies it with that segment.  This holds for every $z$.  Arc length along the segment is precisely the coordinate $d_\psi(x,z)$, proving~\eqref{eq:saturated-line}.
\end{proof}

This applies in particular to the higher-order points used by Alman--Li--Pratt to characterize asymptotic multilinear commutative rank.  Their Section~3.2.2 takes points $\psi$ dominating $\widetilde{\mathrm{CR}}_{1,2}$ on every tensor.  On the $12$ wire of dimension $q$, this rank is $q$.  Hence $\psi(\langle q\rangle_{12})\geq q$, while normalization and tensor-rank monotonicity give the reverse inequality.  Proposition~\ref{prop:saturated-wire-line} therefore makes all their wire metrics line metrics.  Theorem~3.22 of that paper still proves that some of these full tensor functionals lie outside the grouped quantum family; a new full functional need not give a new class of wire metrics.  This conclusion is conditioned on the price-space theorem, not part of the cited rank theorem.

Finally, if $e=c\,d_\psi$ or $e'=c\,d_\psi$, the bound $d_\psi\leq1$ gives $c\geq\sqrt2$.  Equality would saturate a diagonal wire and force a line metric by Proposition~\ref{prop:saturated-wire-line}, which neither configuration is.  Thus $c>\sqrt2$.  This is a restriction on normalized scale, and does not exclude either metric ray from the price cone.

\subsection{Preserving the full spectral point while adjoining midpoints}
The extension statement in the Anthropic preprint can be strengthened from preservation of one tensor to preservation of every old tensor.  The strengthening is existential: it selects suitable extensions, rather than imposing a property on arbitrary spectral extensions.

\begin{proposition}[Normalized extension with all old values fixed]
\label{prop:full-normalized-extension}
Assume Proposition~A.7 of~\cite[page 21]{Anthropic2026}.  Let $I\subseteq I'$ be finite party sets with $|I|\geq2$, let $\pi_j:I'\to I$ be finitely many retractions, and let $\theta_j\geq0$ sum to one.  If $\psi$ is an ordinary normalized spectral point on $I$, there is an ordinary normalized spectral point $\psi'$ on $I'$ such that
\begin{align}
 \psi'(T)&=\psi(T)&&\text{for every tensor supported on }I,\\
 \psi'(T')&\leq\prod_j\psi(\pi_jT')^{\theta_j}
     &&\text{for every }I'\text{-tensor }T'.
 \label{eq:full-normalized-extension}
\end{align}
Here $\pi_jT'$ groups the tensor factors over each fibre of $\pi_j$, and an old tensor is embedded by adjoining one-dimensional factors.
\end{proposition}
\begin{proof}
Fix finitely many nonzero old tensors $T_1,\ldots,T_N$ and apply Proposition~A.7 to
\[
 C=U_2(I)\prod_{a=1}^N T_a.
\]
It gives a relaxed point $\phi'$ bounded above by the product in~\eqref{eq:full-normalized-extension}, with $\phi'(C)=\psi(C)$.  On each old tensor $T$ the upper bound is $\psi(T)$, since every retraction fixes $I$.  Each factor in $C$ has positive spectral value.  Multiplicativity, equality on $C$, and the individual upper bounds therefore force equality on $U_2(I)$ and every $T_a$.

In particular $\phi'(U_2(I))=2$.  Since $U_2(I)$ is a restriction of $U_2(I')$, whose tensor rank is two, we have
\[
 2\leq\phi'(U_2(I'))\leq2.
\]
The normalization exponent on $I'$ is one.  Theorem~\ref{thm:normalization} consequently makes $\phi'$ an ordinary normalized point.

The normalized spectrum on $I'$ is compact in the topology of pointwise convergence, as in the proof of Corollary~\ref{cor:price-attained}.  The displayed upper bounds and the equalities on any specified old tensors are closed conditions.  The preceding construction shows the finite intersection property for the conditions fixing all old values.  Compactness supplies a single normalized point satisfying all of them.  Zero tensors require no additional condition, since every spectral point takes value zero there.
\end{proof}

\begin{theorem}[A coherent normalized point on the whole square]
\label{thm:coherent-square-spectrum}
Assume Proposition~\ref{prop:full-normalized-extension} and the price-space and flat-rectangle results.  Put $K=[0,1]^2$ and let $V$ be its four corners.  Suppose a normalized point $\psi_V$ has $d_{\psi_V}(x,y)=\lambda\|x-y\|$ on $V$, with $\lambda>0$.  There is a family of ordinary normalized points $\psi_F$, one for every finite $V\subseteq F\subseteq K$, such that
\[
 d_{\psi_F}(x,y)=\lambda\|x-y\|\quad(x,y\in F),
 \qquad
 \psi_F(T)=\psi_E(T)\quad(V\subseteq E\subseteq F)
\]
for every tensor $T$ supported on $E$.  The given point on $V$ is preserved.  In particular the centre can be adjoined at the same normalized scale $\lambda$.
\end{theorem}
\begin{proof}
First suppose a finite $F$ with the indicated metric and point is already given, and adjoin the midpoint $m$ of two distinct old points $x,y$.  If $m$ is already present there is nothing to do.  Otherwise apply Proposition~\ref{prop:full-normalized-extension} to the two retractions sending $m$ to $x$ or $y$, each with weight $1/2$.  The new wire distances satisfy
\[
 d(m,x),d(m,y)\leq\tfrac12d(x,y).
\]
The triangle inequality forces equality, so $m$ is the metric midpoint of $x,y$.

We justify that its distances to every old point are the Euclidean ones.  Realize the extended price in a CAT$(0)$ space.  The four corners bound an isometrically embedded Euclidean square of side $\lambda$ by the flat-rectangle result.  Every old point is pinned to its specified position in this square by its four corner distances.  Indeed, for $p=(s,t)\in K$ put
\[
 (w_1,w_2,w_3,w_4)=((1-s)(1-t),s(1-t),st,(1-s)t)
\]
in cyclic corner order.  Twice applying the CAT$(0)$ squared-distance convexity inequality, first along horizontal sides and then along the vertical segment through $p$, gives for any ambient point $u$
\[
 d(u,p)^2\leq\sum_{i=1}^4w_i d(u,v_i)^2
                  -\lambda^2\bigl(s(1-s)+t(1-t)\bigr).
\]
If the corner distances of $u$ are those of the Euclidean point $p$, the right side is zero by direct expansion.  Thus $u=p$.  The old endpoints $x,y$ therefore lie in that flat square in their prescribed positions, and uniqueness of the CAT$(0)$ midpoint puts $m$ at their Euclidean midpoint.  All new distances are exact, while every old tensor value is preserved.

Iterating this construction over successive finite dyadic grids gives compatible normalized points on all such grids.  Restrictions to finite subsets containing $V$ remain ordinary normalized: the old unit tensor $U_2(V)$ has value two, and it is a restriction of the unit tensor on any larger subset.  We obtain a coherent family on every finite subset of the dyadic points containing $V$.

To pass to arbitrary points of $K$, fix a finite set $H\supseteq V$.  Choose injective maps from $H$ into the dyadic points that fix $H$'s dyadic points and converge pointwise to the identity.  Pulling back the already constructed points by these relabelings gives normalized points on $H$.  Compactness supplies a limit with the required Euclidean wire values and with all tensor values on $H$'s dyadic subset preserved.  Finally take the product of the compact normalized spectra over all finite $H\supseteq V$.  Impose the closed wire-value and restriction-consistency conditions.  Any finite collection of them is satisfied by constructing a point on the union of the finitely many label sets involved and restricting it.  The finite intersection property therefore supplies one coherent family for all finite subsets of $K$.
\end{proof}

Combining this theorem with Proposition~\ref{prop:collinear-edge-point} fixes the full spectral functional on every collinear triple in the square, with one common scale.  If $z$ lies between $x$ and $y$, then
\[
 \psi_F(T)=\left(\zeta^{(\rho,1-\rho,0)}(T)\right)^{\lambda\|x-y\|},
 \qquad \rho=\frac{\|x-z\|}{\|x-y\|},
\]
for every tensor supported on that triple, with slots ordered $(x,y,z)$.  Consequently an identity involving products of such line-supported tensors imposes a computable constraint on the same spectral point.  For example, if $P_j$ is a product of such tensors and
\[
 L_j=\sum_{T\text{ a factor of }P_j}\|x_T-y_T\|
       \log_2\zeta^{(\rho_T,1-\rho_T,0)}(T),
\]
then a cover of $U_{q_0}(F)P_0$ by restrictions of the tensors $U_{q_j}(F)P_j$ forces
\[
 q_0\,2^{\lambda L_0}\leq\sum_j q_j\,2^{\lambda L_j}.
\]
The theorem selects one coherent extension preserving every old tensor value.  Corollary~\ref{cor:unit-geodesic-insertion} gives a further statement for arbitrary normalized points: adjoining a party at a prescribed point of a geodesic between two old parties does not increase the unit exponent of a subset containing those endpoints.  Thus this equality holds for every extension with that geodesic placement, although arbitrary placements are outside its scope.  No contradiction among the resulting mixed-line constraints has been established here.

\subsection{The crossing projectors have a communication cost}

The nonzero projector products above cannot be inserted as free local
operations in a wire-network identity. The following exact calculation
quantifies that issue on the very sector used in the witness.

\begin{proposition}\label{prop:crossing-projector-cost}
On the two-dimensional local $[22]$ module, the operator Schmidt ranks
of the antisymmetric-line projector and of the symmetric-traceless
projector on the tensor square are respectively $4$ and $2$.
If the two input and output ports are located at points $x,y$ of a
Euclidean space, the least cable cost of a single exact implementation
of these entire operators is respectively $2\|x-y\|$ and $\|x-y\|$.
For tensor powers of these fixed operators, both costs multiply by the
number of powers.
\end{proposition}

\begin{proof}
Choose an orthonormal basis of the module and put
$J=\left(\begin{smallmatrix}0&-1\\1&0\end{smallmatrix}\right)$.
The symmetric-traceless projector is
$(I\otimes I-J\otimes J)/2$, which has operator Schmidt rank two
because $I,J$ are linearly independent. The antisymmetric projector
is $|\varepsilon\rangle\langle\varepsilon|/2$, where
$\varepsilon=e_1\otimes e_2-e_2\otimes e_1$. Its coefficient matrix
across the two sites, after grouping each input with its output, is
the tensor product of two invertible matrices, up to permutations
and a nonzero scalar. Its rank is therefore four.

Vectorization identifies a bipartite operator of operator Schmidt
rank $r$ with a two-party tensor of matrix rank $r$. A dimension-$r$
wire realizes it by local linear maps. Conversely the cut-rank
integration in Proposition~\ref{prop:single-layout} bounds every
single realization, including ones with extra hubs, by
$\log_2(r)\|x-y\|$. Matrix rank multiplies under tensor product,
which proves the tensor-power assertion.
\end{proof}

This bound is for implementing the whole operator on arbitrary inputs.
Its value on one special input may be cheaper. A sum of implementations
or a different collective projector on a larger copy block is also
outside the statement. Thus the calculation explains a cost that must
be addressed before the upper-functional counterexample could become
a universal spectral inequality; it is not itself a square nonprice
certificate.
